% !TeX root = ../../Book3.tex
% 纲要标题：### H.3 Serre 对偶与上同调计算
% 状态：正文与算例初稿；一般对偶给证明概要，射影直线配对逐项计算。

\section{Serre 对偶与上同调计算}
\label{b3:sec:H03}

在射影直线的两张标准图上，\(t^{-1}\) 代表 \(H^1(\mathcal O(-2))\) 中不能消去的类。把它乘以微分 \(\mathrm{d}t\) 后，取 Laurent 展开中 \(t^{-1}\mathrm{d}t\) 的系数，就得到一个数。一般曲线上的 Serre 对偶把这种计算组织为一次上同调与微分截面之间的完美配对；\secref{b3:sec:H02}所用的修正项 \(\ell(K_C-D)\) 正由此出现。

\ProseParagraphBreak
以下仍设 \(k\) 为代数闭域，\(C\) 为光滑射影整曲线。记 \(\omega_C=\Omega_{C/k}\)，并以 \(V^\vee=\operatorname{Hom}_k(V,k)\) 表示有限维 \(k\) 空间的对偶。一般配对的存在性先用射影嵌入给出证明概要；两图上的配对则直接由截面复形算出。

\subsection{曲线上的对偶配对}
\label{b3:subsec:H03:01}

若 \(\mathcal L\) 是线丛，一个截面 \(s\in H^0(C,\omega_C\otimes\mathcal L^{-1})\) 等价于层同态 \(s:\mathcal L\rightarrow\omega_C\)。它在一次上同调上诱导线性映射。因而只要给出 \(H^1(C,\omega_C)\) 到 \(k\) 的适当迹映射，就可把一次上同调类与 \(s\) 配成一个数。

\begin{theorem}[曲线上的 Serre 对偶]\label{b3:thm:H-serre-duality}
	设 \(k\) 为代数闭域，\(C\) 为光滑射影整曲线。存在迹同构
	\[
		\operatorname{Tr}_C:H^1(C,\omega_C)\longrightarrow k,
		\qquad\omega_C=\Omega_{C/k},
	\]
	使对每个线丛 \(\mathcal L\)，下列配对都是有限维 \(k\) 空间之间的完美配对：
	\[
		\begin{aligned}
			H^1(C,\mathcal L)\times H^0(C,\omega_C\otimes\mathcal L^{-1})
			&\longrightarrow k,\\
			(\alpha,s)&\longmapsto\operatorname{Tr}_C(H^1(s)(\alpha)).
		\end{aligned}
	\]
	换言之，这一配对诱导自然同构
	\[
		H^1(C,\mathcal L)\cong H^0(C,\omega_C\otimes\mathcal L^{-1})^\vee.
	\]
	把 \(\mathcal L\) 换成 \(\omega_C\otimes\mathcal L^{-1}\)，也得到 \(H^0(C,\mathcal L)\) 与 \(H^1(C,\omega_C\otimes\mathcal L^{-1})\) 的完美配对。
\end{theorem}
\begin{proofsketch}
	取闭嵌入 \(i:C\hookrightarrow P=\mathbb P_k^N\)，记 \(c=N-1\)、\(\omega_P=\bigwedge^N\Omega_{P/k}\)。射影空间上的对偶首先给出
	\[
		\operatorname{Ext}^{c}_{\mathcal O_P}(i_*\mathcal L,\omega_P)
		\cong H^1(P,i_*\mathcal L)^\vee
		\cong H^1(C,\mathcal L)^\vee.
	\]
	第一步是\cite[Chapter III, Theorem 7.1, p. 240]{Hartshorne1977}：先对 \(\mathcal O_P(a)\) 用标准仿射覆盖上的单项式配对，再对有限个扭曲层的直和作同样计算；一般凝聚层由这些直和展示，零次的 Hom 配对由正合列延拓，高次 Ext 则由可消去性和导出函子的比较得到。第二步可在 \(P\) 的仿射覆盖上检验：闭嵌入的逆像仍仿射，\(i_*\mathcal L\) 的截面与 \(\mathcal L\) 的截面相同，两份 Čech 复形逐项一致。

	为把左端写回 \(C\) 上，设 \(\mathcal I\) 为 \(C\) 的理想层，取
	\[
		\omega'_C\coloneqq
		\mathcal{E}xt^c_{\mathcal O_P}(i_*\mathcal O_C,\omega_P),
	\]
	把它视作 \(C\) 上的层。在光滑嵌入的局部，\(\mathcal I\) 由长度 \(c\) 的正则序列生成。该序列的 Koszul 复形是 \(\mathcal O_C\) 的局部自由消解，其 Hom 复形在 \(c\) 次以前无上同调。截取内射消解时，这些较低次的正合项可以消去，得到
	\[
		\operatorname{Hom}_C(\mathcal L,\omega'_C)
		\cong\operatorname{Ext}^c_{\mathcal O_P}(i_*\mathcal L,\omega_P).
	\]
	这一 Hom 与 Ext 的变换，以及由它得到的对偶化层构造，分别见\cite[Chapter III, Lemma 7.4 and Proposition 7.5, p. 242]{Hartshorne1977}。其技术内容包括在层范畴中保持内射性的伴随识别，以及低次正合复形的分裂截取。

	同一 Koszul 计算还说明，更换正则序列生成元时，最高外幂带来行列式因子；用余法丛的行列式对偶补偿后，局部同构黏合为
	\[
		\omega'_C\cong\omega_P|_C\otimes\det(\mathcal I/\mathcal I^2)^\vee
		\cong\Omega_{C/k}.
	\]
	最后一个同构来自光滑闭嵌入的余法正合列
	\[
		0\longrightarrow\mathcal I/\mathcal I^2
		\longrightarrow\Omega_{P/k}|_C
		\longrightarrow\Omega_{C/k}\longrightarrow0
	\]
	的最高外幂。这里采用\cite[Chapter III, Theorem 7.11 and Corollary 7.12, pp. 245--246]{Hartshorne1977}的局部完全交计算及典范层识别。

	合成上述同构，得到对 \(\mathcal L\) 自然的
	\[
		\operatorname{Hom}_C(\mathcal L,\omega_C)
		\cong H^1(C,\mathcal L)^\vee.
	\]
	令 \(\mathcal L=\omega_C\)，恒等同态对应一个泛函 \(\operatorname{Tr}_C\)。由自然性，任意 \(s:\mathcal L\rightarrow\omega_C\) 所对应的泛函正是 \(\operatorname{Tr}_C\circ H^1(s)\)，这给出定理中的配对。又因 \(\operatorname{Hom}_C(\omega_C,\omega_C)=H^0(C,\mathcal O_C)=k\)，迹映射为同构。

	这里省略射影空间上一般凝聚层的导出函子比较、光滑嵌入的正则序列构造，以及层 Ext 的截取和黏合细节；上述文献给出了这些步骤的证明。一般局部自由层的对偶形式另见\cite[Chapter III, Corollary 7.7, p. 244; Corollary 7.12 and Remark 7.12.2, p. 246]{Hartshorne1977}。
\end{proofsketch}

这一定理称为\term[Serre-duiou]{Serre 对偶}[Serre duality]。对偶空间中的微分带有 \(\mathcal L^{-1}\)；特别地，当 \(\mathcal L=\mathcal O_C(D)\) 时，它是 \(\omega_C\otimes\mathcal O_C(-D)\cong\mathcal O_C(K_C-D)\)，除子的符号由此确定。

\ProseParagraphBreak
配对还可直接传到 Čech 代表元。对有限仿射开覆盖 \((U_i)\)，取表示 \(\alpha\) 的交叠截面 \(a_{ij}\in\mathcal L(U_i\cap U_j)\)。乘以全局截面 \(s\) 后，\((s a_{ij})\) 是 \(\omega_C\) 的余圈。乘法与限制相容，且把余边送到余边，因此所得上同调类只依赖 \(\alpha,s\)。\cref{b3:thm:affine-cech-comparison}识别这份计算与层上同调；迹随后将它送到 \(k\)。

\ProseParagraphBreak
若 \(u:\mathcal L\rightarrow\mathcal M\) 是线丛同态，则 \(s\circ u\) 与 \(H^1(u)(\alpha)\) 满足
\[
	\langle H^1(u)(\alpha),s\rangle_{\mathcal M}
	=\langle\alpha,s\circ u\rangle_{\mathcal L}.
\]
因而在完美配对的同构下，有交换图
\[
	\begin{tikzcd}[column sep=large]
		H^1(C,\mathcal L)\arrow[r,"H^1(u)"]\arrow[d,"\cong"']
		&H^1(C,\mathcal M)\arrow[d,"\cong"]\\
		H^0(C,\omega_C\otimes\mathcal L^{-1})^\vee
		\arrow[r,"(u^*)^\vee"']
		&H^0(C,\omega_C\otimes\mathcal M^{-1})^\vee .
	\end{tikzcd}
\]
这里 \(u^*\) 表示把到 \(\omega_C\) 的同态与 \(u\) 复合。这个方向说明，上同调映射的对偶是微分截面上的反向映射。

\subsection{线丛上同调与典型计算}
\label{b3:subsec:H03:02}

射影直线上的迹可以明确指定，不必保留为抽象泛函。取
\[
	U_0=\operatorname{Spec}k[t],\qquad
	U_\infty=\operatorname{Spec}k[u],\qquad u=t^{-1}.
\]
对 \(\mathcal O(n)=\mathcal O(nP_\infty)\)，以 \(e_0=1\)、\(e_\infty=t^n\) 为局部基，所以 \(e_\infty=t^ne_0\)。采用“后一张减前一张”的 Čech 微分，截面复形为
\[
	k[t]e_0\oplus k[t^{-1}]e_\infty
	\longrightarrow k[t,t^{-1}]e_0,\qquad
	(ae_0,be_\infty)\longmapsto(t^n b-a)e_0.
\]
两图及交集都仿射，\cref{b3:thm:affine-cech-comparison}保证它计算层上同调。核由 \(a(t)=t^nb(t^{-1})\) 给出；第一边的指数非负，第二边的指数至多为 \(n\)，所以核的基为 \(e_0,te_0,\ldots,t^ne_0\)（\(n\ge0\)），\(n<0\) 时为零。余核则为
\[
	H^1(\mathbb P^1,\mathcal O(n))
	\cong\frac{k[t,t^{-1}]}{k[t]+t^nk[t^{-1}]}\,e_0.
\]
分母消去指数 \(j\ge0\) 和 \(j\le n\)，留下的基为 \([t^je_0]\)、\(n<j<0\)。这同时给出
\[
	h^0(\mathcal O(n))=\max\{n+1,0\},\qquad
	h^1(\mathcal O(n))=\max\{-n-1,0\},
\]
并与\cref{b3:prop:projective-line-cohomology}的计算一致。

\begin{proposition}[射影直线上的显式配对]\label{b3:prop:H-projective-line-duality}
	设 \(k\) 为代数闭域，\(n\in\mathbb Z\)。在上述坐标与局部基下，取迹同构
	\[
		\operatorname{Tr}_{\mathbb P^1}([h(t)\,\mathrm{d}t])
		=\operatorname{coeff}_{t^{-1}}h(t),
	\]
	其中右端表示 Laurent 多项式中 \(t^{-1}\) 的系数。若 \(n\le-2\)，则
	\[
		\begin{aligned}
			H^1(\mathcal O(n))&=\bigoplus_{i=1}^{-n-1}k\cdot[t^{-i}e_0],\\
			H^0(\omega_{\mathbb P^1}\otimes\mathcal O(-n))
			&=\bigoplus_{j=0}^{-n-2}k\,t^j\mathrm{d}t\otimes e_0^{-1},
		\end{aligned}
	\]
	配对为
	\[
		\langle[he_0],q\,\mathrm{d}t\otimes e_0^{-1}\rangle
		=\operatorname{coeff}_{t^{-1}}(hq).
	\]
	特别地，\([t^{-i}e_0]\) 与 \(t^{i-1}\mathrm{d}t\otimes e_0^{-1}\) 互为对偶基。若 \(n\ge-1\)，这两个空间都为零。
\end{proposition}
\begin{proof}
	在第二张图上 \(\mathrm{d}u=-t^{-2}\mathrm{d}t\)，所以正则微分的两图贡献为 \(k[t]\mathrm{d}t\) 与 \(t^{-2}k[t^{-1}]\mathrm{d}t\)。它们都没有 \(t^{-1}\mathrm{d}t\) 项，而商只剩这一项。因此所写迹良定义且为同构。

	对 \(\omega_{\mathbb P^1}\otimes\mathcal O(-n)\)，局部基分别为 \(\mathrm{d}t\otimes e_0^{-1}\) 和 \(\mathrm{d}u\otimes e_\infty^{-1}\)，后者等于前者的 \(-t^{-n-2}\) 倍。全局截面因而写成 \(q(t)\mathrm{d}t\otimes e_0^{-1}\)，其中 \(q\in k[t]\)、\(\deg q\le-n-2\)。若这个上界为负，只有零截面。

	乘积与迹给出所列公式。若把 \(h\) 增加一个 \(a(t)\in k[t]\)，则 \(a q\) 的指数非负，不影响 \(t^{-1}\) 系数；若增加 \(t^n b(t^{-1})\)，乘 \(q\) 后的最高指数至多为 \(n+(-n-2)=-2\)，也不影响该系数。因此配对在所写余核上良定义。最后，对 \(1\le i,r\le-n-1\)，\(t^{-i}t^{r-1}\) 的 \(t^{-1}\) 系数在 \(i=r\) 时为一，否则为零，故为完美配对。\(n\ge-1\) 时两端都为零，结论同样成立。
\end{proof}

在另一方向，当 \(n\ge0\) 时，\(H^0(\mathcal O(n))\) 的基 \(t^ie_0\)（\(0\le i\le n\)）与
\[
	[t^{-i-1}\mathrm{d}t\otimes e_0^{-1}]
	\in H^1(\omega_{\mathbb P^1}\otimes\mathcal O(-n))
\]
互为对偶基。这里的相同系数计算解释了为什么零次截面的非负指数，与另一线丛一次上同调的负指数恰好对应。上述迹的标准 Čech 生成元也见\cite[Chapter III, Remark 7.1.1, pp. 240--241]{Hartshorne1977}。

\ProseParagraphBreak
回到一般曲线，若有效除子 \(Z\) 的次数为 \(m\)，则限制线丛 \(\mathcal L|_Z\) 的全局截面空间为 \(m\) 维：在每个支撑点，长度为该点系数的局部商环上，线丛仍为秩一自由模。短正合列
\[
	0\longrightarrow\mathcal L(-Z)\longrightarrow\mathcal L
	\longrightarrow\mathcal L|_Z\longrightarrow0
\]
给出求值映射及其障碍：
\[
	\begin{aligned}
		0&\longrightarrow H^0(\mathcal L(-Z))\longrightarrow H^0(\mathcal L)
		\xrightarrow{\mathrm{ev}_Z}H^0(\mathcal L|_Z)\\
		&\xrightarrow{\delta}H^1(\mathcal L(-Z))\longrightarrow H^1(\mathcal L)
		\longrightarrow0.
	\end{aligned}
\]
在 \(Z=2P\) 时，求值指定的是局部展开的常数项和一次项。若 \(\deg\mathcal L-m>2g-2\)，\cref{b3:cor:H-canonical-degree}使 \(H^1(\mathcal L(-Z))=0\)，所以这些局部条件都能同时满足；次数较小时，\(\delta\) 正记录不能满足的条件。Serre 对偶把这一障碍空间识别为微分截面空间的对偶。

\begin{example}[次数为零的两种上同调]\label{b3:exm:H-degree-zero-cohomology}
	设 \(k\) 为代数闭域，\(C\) 为亏格 \(g\ge1\) 的光滑射影整曲线，\(P,Q\) 为不同闭点。线丛 \(\mathcal O_C(P-Q)\) 不平凡：若 \(P-Q\) 为主除子，就存在只有一个简单极点的非恒定函数，由\cref{b3:prop:H-curve-morphism}得到到 \(\mathbb P^1\) 的次数一态射，进而 \(C\cong\mathbb P^1\)，与 \(g\ge1\) 矛盾。由\cref{b3:cor:H-section-degree}，该线丛的零次上同调为零，而 Euler 特征公式给出
	\[
		\begin{aligned}
			\bigl(h^0(\mathcal O_C),h^1(\mathcal O_C)\bigr)&=(1,g),\\
			\bigl(h^0(\mathcal O_C(P-Q)),h^1(\mathcal O_C(P-Q))\bigr)&=(0,g-1).
		\end{aligned}
	\]
	两者次数均为零，上同调却不同。特别地，亏格一曲线上的非平凡零次数线丛没有任何上同调；亏格至少二时，其一次上同调仍非零。
\end{example}

\begin{proposition}\label{b3:prop:H-genus-one-cohomology}
	设 \(k\) 为代数闭域，\(C\) 为亏格一的光滑射影整曲线，\(\mathcal L\) 为次数 \(d\) 的线丛。则
	\[
		\bigl(h^0(C,\mathcal L),h^1(C,\mathcal L)\bigr)=
		\begin{cases}
			(d,0),&d>0,\\
			(1,1),&\mathcal L\cong\mathcal O_C,\\
			(0,0),&d=0\;\text{且}\;\mathcal L\not\cong\mathcal O_C,\\
			(0,-d),&d<0.
		\end{cases}
	\]
\end{proposition}
\begin{proof}
	\cref{b3:cor:H-canonical-degree}给出 \(\omega_C\cong\mathcal O_C\)，所以 Serre 对偶把 \(H^1(C,\mathcal L)\) 识别为 \(H^0(C,\mathcal L^{-1})^\vee\)。任意线丛由\cref{b3:prop:H-divisor-line-bundle}写成 \(\mathcal O_C(D)\)，再对 \(D\) 和 \(-D\) 分别应用\cref{b3:prop:H-genus-one-sections}，便得到所列四种情形。
\end{proof}

这张维数表中，正次数的消失与负次数的非消失互为对偶；零次数还必须区分平凡与非平凡线丛。一般亏格下也有同样的计算方式：先辨认 \(K_C-D\) 的有效截面，再由对偶确定一次上同调，最后用 Euler 特征恢复零次维数。次数界能使某些项直接消失，低次数时则需要具体的除子类信息。


% 新增结构与纲要同步；正文与算例初稿，尚未正式审核。
% #### H.3.1 曲线上的对偶配对
% #### H.3.2 线丛上同调与典型计算
% * 状态：初稿已完成；新增范围尚未正式审核。
